%!TEX root = ../thesis.tex

\chapter{Introduction}\label{chap:introduction}

% Short framing chapter (~5 pages). The substance lives in \cref{chap:background}
% (fuzzing, RL) and \cref{chap:related-work} (RL for fuzzing); this chapter only
% states what the thesis asks, what it contributes, and how it is organized.

\section{Motivation}\label{sec:introduction-motivation}

% TODO: Condense the argument that is developed at length in \cref{chap:background}:
%   - coverage-guided greybox fuzzing is the dominant technique but plateaus
%     (\cref{sec:fuzzing-limitations}),
%   - stateful/reactive targets break its one-input-per-execution model
%     (\cref{sec:fuzzing-stateful}),
%   - emulation restores observability and a cheap exact reset, at the cost of
%     throughput (\cref{sec:fuzzing-emulation}) --- which makes sample efficiency,
%     not executions per second, the binding constraint,
%   - hence a sample-efficient, model-based RL agent (EfficientZero) rather than
%     the model-free learners used by prior RL fuzzers.

\section{Problem Statement and Research Questions}\label{sec:introduction-rqs}

% TODO: State the problem in one paragraph, then the research questions. These are
% currently drafted as prose bullets in \cref{sec:rlfuzz-gap}; promote them here as a
% numbered list and cite them by number from \cref{chap:results} onward, so every
% result answers a named question.
%
% Proposed numbering (each RQ maps to one results section):
%   RQ1  Can EfficientZero learn useful policies on fuzzing-shaped reward
%        landscapes (sparse, decaying, heavily aliased observations)?
%        -> \cref{sec:results-input-actions}
%   RQ2  How should the fuzzing action space be modelled --- as the input the
%        program consumes, or as operations on a seed corpus?  What does each
%        formulation cost and buy?
%        -> \cref{sec:results-head-to-head}
%   RQ3  Does planning with a learned model help, or is the benefit attributable
%        to the learned value function alone?
%        -> \cref{sec:results-ablations} (sims=1 vs.\ full MCTS)
%   RQ4  Does the approach transfer from diagnostic microbenchmarks
%        (\cref{sec:targets-microbenchmarks}) to targets whose reward was not
%        designed for it (\cref{sec:targets-realistic})?
%        -> \cref{sec:results-corpus-actions,sec:discussion-threats}

\section{Contributions}\label{sec:introduction-contributions}

% TODO: Enumerate, one sentence each:
%   1. Two MDP formulations of emulator-based greybox fuzzing that satisfy the
%      planning contract of a MuZero-family learner (\cref{chap:fuzzing-mdp}),
%      and a controlled comparison of the two on a shared target
%      (\cref{sec:results-head-to-head}).
%   2. An implementation: a reusable EfficientZero in Rust on burn, and its
%      integration with the MOGI emulator as a fuzzing driver
%      (\cref{chap:implementation}).
%   3. A diagnostic methodology for model-based RL collapse in this setting ---
%      probes, failure taxonomy, and the evidence that distinguishes the modes
%      (\cref{chap:failure-modes}).
%   4. An empirical evaluation with a pre-registered multi-seed protocol,
%      reporting both executions and wall-clock (\cref{chap:methodology,chap:results}).

\section{Outline}\label{sec:introduction-outline}

% TODO: One short paragraph walking the reader through the parts:
%   background (\cref{chap:background,chap:efficientzero,chap:related-work}),
%   approach (\cref{chap:platform,chap:fuzzing-mdp,chap:input-actions,chap:corpus-actions,chap:implementation}),
%   evaluation (\cref{chap:targets,chap:methodology,chap:results,chap:failure-modes,chap:discussion}).
